\exercisesection{Macaulay 模与环}

\begin{enumerate}
\def\labelenumi{\arabic{enumi})}
\item
  设 A 为 Noether 局部环，并设 \(0 \to M' \to M \to M'' \to 0\) 为 Macaulay A-模的一个正合列。证明我们有 \(\dim(M') = \dim(M)\) ，且 \(\dim(M'')\) 等于 \(\dim(M)\) 或 \(\dim(M) - 1\) 。
\item
  设 A 为 Macaulay 局部环，并设 \(\mathfrak{p} \in \operatorname{Spec}(A)\) 。证明我们有 \(\dim(\widehat{A}/\mathfrak{q}) = \dim(A/\mathfrak{p})\) （对每个 \(\mathfrak{q} \in \operatorname{Ass}(\widehat{A}/\mathfrak{p}\widehat{A})\) ）。特别地，环 \(\widehat{\mathrm{A}}/\mathfrak{p}\widehat{A}\) 没有嵌入素理想。
\item
  \begin{enumerate}
  \def\labelenumii{\alph{enumii})}
  \tightlist
  \item
    设 R 为一个 Noether 的、整的、整闭的 Q-代数，A 为环，并设 \(\rho: \mathrm{R} \to \mathrm{A}\) 为使 A 成为有限 R-代数的一个单射同态。证明子 R-模 \(\rho(\mathrm{R})\) 是 A 的一个直和项（化为命题 8 的推论 3）。
  \end{enumerate}
\end{enumerate}

\begin{enumerate}
\def\labelenumi{\alph{enumi})}
\setcounter{enumi}{1}
\tightlist
\item
  设 A 为一个 Noether 局部 Q-代数。证明 A 满足下列性质（“单项式条件”）：
\end{enumerate}

\begin{enumerate}
\def\labelenumi{(\Roman{enumi})}
\setcounter{enumi}{899}
\tightlist
\item
  对每个极大割序列 \((x_{1},\ldots,x_{d})\) （由 \(m_{A}\) 中的元素组成）以及每个整数 p，元素 \((x_{1}\ldots x_{d})^{p}\) 不属于理想 \((x_{1}^{p+1},\ldots,x_{d}^{p+1})\) 。（化为 A 是完备的情形，从而它具有一个系数域 K，并对一个同态 \(\rho:K[[X_{1},\ldots,X_{d}]]\to A\) （它满足 \(\rho(X_{i})=x_{i}\) ）应用 a）。）
\end{enumerate}

\begin{enumerate}
\def\labelenumi{\arabic{enumi})}
\setcounter{enumi}{3}
\tightlist
\item
  设 A 为局部环，p 为 A 的一个素理想。用 B 表示分式环 \(A[Z][S^{-1}]\) ，其中 S 是形如 \(Z + a\) 的元素（在 A{[}Z{]}中）组成的集合（\(a \in \mathfrak{m}_{A} - \mathfrak{p}\) ）。用 z 表示 B 中的元素 Z/1。
\end{enumerate}

\begin{enumerate}
\def\labelenumi{\alph{enumi})}
\item
  证明 \(\mathrm{B} / (z)\) 同构于 \(\mathbf{A}_{\mathfrak{p}}\) ，且 \(\mathrm{B} / (z - 1)\) 同构于 \(\mathbf{A}\) 。
\item
  以下假设环 A 是整的且 Noether 的，并假设 p 由 A 的一个非零元素 p 生成。设 C 为环 B{[}T{]}/(z(pT-1))；若 A 是 Macaulay 环，则 C 也是。证明 Spec C 有两个不可约分支 \(X_{1}\) 与 \(X_{2}\) ，它们相交于一个闭点 x，使得 \(\text{codim}(\{x\}, X_{1}) = \text{codim}(\{x\}, X_{2}) = 1\) ，\(\text{dim}(X_{1}) = 2\) 且 \(\text{dim}(X_{2}) \geqslant \text{dim}(A_{\mathfrak{p}})\) 。对每个整数 n，给出一个满足 \(\text{dim}(A_{\mathfrak{p}}) = n\) 的例子。
\end{enumerate}

¶ 5) 设 A 为 Noether 局部环，M 为有限生成 A-模，\(\mathbf{x} = (x_{1},\ldots ,x_{n})\) 为 \(\mathfrak{m}_{\mathrm{A}}\) 中元素组成的一个序列，它生成一个理想 \(\mathfrak{x}\)；假设 A-模 M/xM 有有限长度（这蕴含 \(n\geqslant \dim (\mathrm{M})\) ）。用 \(\mathrm{gr(A)}\) 与 \(\mathrm{gr(M)}\) 表示 A 与 M 关于 x-adic 拓扑的伴随分次环，用 \(\xi_{i}\) （ \(1\leqslant i\leqslant n\) ）表示 \(x_{i}\) 在 \(\mathfrak{x} / \mathfrak{x}^{2}\) 中的像，用 \(\xi\) 表示序列 \((\xi_1,\dots ,\xi_n)\) 。

\begin{enumerate}
\def\labelenumi{\alph{enumi})}
\item
  对 \(p \in \mathbf{Z}\) ，令 \(\mathrm{F}^p\mathbf{K}_i(\mathbf{x},\mathrm{M}) = \mathbf{K}_i(\mathbf{x},\mathfrak{x}^{p - i}\mathrm{M})\) （约定 \(\mathfrak{x}^s = A\) ，对 \(s \leqslant 0\) ）。证明 \(\mathrm{F}^p\mathbf{K}_{\bullet}(\mathbf{x},\mathrm{M})\) 是 \(\mathbf{K}_{\bullet}(\mathbf{x},\mathrm{M})\) 的一个子复形，且复形 \(\bigoplus_{n}\mathrm{F}^p\mathbf{K}_{\bullet}(\mathbf{x},\mathrm{M}) / \mathrm{F}^{p + 1}\mathbf{K}_{\bullet}(\mathbf{x},\mathrm{M})\) 等同于 \(\mathbf{K}_{\bullet}^{\mathrm{gr}(A)}(\boldsymbol {\xi},\mathrm{gr}(M))\) 。
\item
  对每个使 H(C) 具有有限长度的 A-模复形 C，令 \(\chi(C) = \sum_{i}(-1)^{i}\lg(\mathrm{H}_{i}(C))\) 。证明存在一个整数 \(p_0\) ，使得复形 \(\mathrm{F}^{p}\mathbf{K}_{\bullet}(\mathbf{x},\mathrm{M}) / \mathrm{F}^{p + 1}\mathbf{K}_{\bullet}(\mathbf{x},\mathrm{M})\) 的同调为零（对 \(p\geqslant p_0\) ）；由此导出等式
\end{enumerate}

\[
\chi \left(\mathbf {K} _ {\bullet} ^ {\operatorname{gr} (\mathrm{A})} (\boldsymbol {\xi}, \operatorname{gr} (\mathrm{M}))\right) = \chi \left(\mathbf {K} _ {\bullet} (\mathrm{x}, \mathrm{M}) / \mathrm{F} ^ {p} \mathbf {K} _ {\bullet} (\mathrm{x}, \mathrm{M})\right) = \sum (- 1) ^ {i} \binom {n} {i} \lg (\mathrm{M} / x ^ {p - i} \mathrm{M})
\]

对 \(p \geqslant p_0\) 成立。

\begin{enumerate}
\def\labelenumi{\alph{enumi})}
\setcounter{enumi}{2}
\item
  证明上列表达式当 \(n > \dim(\mathrm{M})\) 时为零，且等于 \(e_{\mathfrak{x}}(\mathrm{M})\) （当 \(n = \dim(\mathrm{M})\) 时）。
\item
  设 \(h \geqslant 0\) 且 \(p \geqslant p_{0}\)；证明从 \(\mathrm{F}^{p}\mathbf{K}_{\bullet}(\mathbf{x},\mathrm{M})\) 到 \(\mathrm{F}^{p+h}\mathbf{K}_{\bullet}(\mathbf{x},\mathrm{M})\) 的典范单射是一个同调同构。由此推出 \(\mathrm{H}_{i}(\mathrm{F}^{p}\mathbf{K}_{\bullet}(\mathbf{x},\mathrm{M}))\) 对 x-adic 拓扑是离散的。进而证明这个模的由映射 \(\mathrm{H}_{i}(\mathrm{F}^{p+h}\mathbf{K}_{\bullet}(\mathbf{x},\mathrm{M})) \to \mathrm{H}_{i}(\mathrm{F}^{p}\mathbf{K}_{\bullet}(\mathbf{x},\mathrm{M}))\) 的像所定义的滤链是 x-好的，并由此推出 \(\mathrm{F}^{p}\mathbf{K}_{\bullet}(\mathbf{x},\mathrm{M})\) 的同调为零。最后推出 \(\chi(\mathbf{K}_{\bullet}(\mathbf{x},\mathrm{M}))\) 当 \(n > \dim(\mathrm{M})\) 时为零，且等于 \(e_{\mathfrak{x}}(\mathrm{M})\) （当 \(n = \dim(\mathrm{M})\) 时）。
\end{enumerate}

\begin{enumerate}
\def\labelenumi{\arabic{enumi})}
\setcounter{enumi}{5}
\tightlist
\item
  我们保留上一习题的假设；对每个使 N/xN 具有有限长度的有限生成 A-模 N，用 \(\chi(\mathbf{x},\mathbf{N})\) 表示整数 \(\chi(\mathbf{K}_{\bullet}(\mathbf{x},\mathbf{N}))\) 。
\end{enumerate}

\begin{enumerate}
\def\labelenumi{\alph{enumi})}
\item
  对每个子模 \(M'\) （属于 \(M\) ），有 \(\chi(\mathbf{x}, M) = \chi(\mathbf{x}, M') + \chi(\mathbf{x}, M/M')\) 。
\item
  若 \(x_{1}\mathrm{M} = 0\)，证明 \(\chi (\mathbf{x},\mathrm{M})\) 为零。
\item
  用 \(\mathbf{x}'\) 表示序列 \((x_2,\ldots ,x_n)\)；证明等式
\end{enumerate}

\[
\chi (\mathbf {x}, \mathrm{M}) = \chi (\mathbf {x} ^ {\prime}, \mathrm{M} / x _ {1} \mathrm{M}) - \chi (\mathbf {x} ^ {\prime}, \operatorname{Ker} (x _ {1}) _ {\mathrm{M}})
\]

（利用 A, X, 第 207 页，习题 8）。

\begin{enumerate}
\def\labelenumi{\alph{enumi})}
\setcounter{enumi}{3}
\tightlist
\item
  对 \(1 \leqslant i \leqslant n\) ，用 \(\mathrm{K}_i\) 表示以 \(x_i\) 为比的乘法映射（在 \(\mathrm{M}/(x_1\mathrm{M} + \ldots + x_{i-1}\mathrm{M})\) 中的）的核。证明等式
\end{enumerate}

\[
\lg (\mathrm{M} / \mathfrak {x} \mathrm{M}) - \chi (\mathbf {x}, \mathrm{M}) = \sum_ {i = 1} ^ {n} \chi ((x _ {i + 1}, \dots , x _ {n}), \mathrm{K} _ {i}).
\]

\begin{enumerate}
\def\labelenumi{\alph{enumi})}
\setcounter{enumi}{4}
\tightlist
\item
  设 \(p_1, \ldots, p_n\) 为 \(> 1\) 的整数，\(\mathbf{x}^{\mathbf{p}}\) 为序列 \((x_{1}^{p_1}, \ldots, x_n^{p_n})\) 。证明等式 \(\chi(\mathbf{x}^{\mathbf{p}}, \mathrm{M}) = p_1 \ldots p_n \chi(\mathbf{x}, \mathrm{M})\) （先处理 \(n = 1\) 的情形，然后利用 A, X, 第 157 页，推论 2 化为这一情形）。
\end{enumerate}

¶ 7) 设 A 为 Noether 局部环，M 为有限生成 A-模。我们称序列 \((x_1, \ldots, x_n)\) 对 M 是弱正则的，如果对 \(i = 1, \ldots, n\) ，以 \(x_{i+1}\) 为比的乘法映射（在 \(\mathrm{M}/(x_1\mathrm{M} + \ldots + x_i\mathrm{M})\) 中的）的核被 \(\mathfrak{m}_A\) 零化。

\begin{enumerate}
\def\labelenumi{\alph{enumi})}
\item
  若 \(n \leqslant \dim(M)\)，证明对 M 弱正则的序列 \((x_{1}, \ldots, x_{n})\) 对 M 是割的（对 n 用归纳法推理）。给出一个弱正则但不割的序列的例子（可取 \(\mathrm{A} = M = k[X]/(X^{2})\) （其中 k 是一个域）。
\item
  我们称 M 为 Buchsbaum 模，如果对 M 割的每个序列都对 M 弱正则。Macaulay 模是 Buchsbaum 模；若 M 是 Buchsbaum 模，则 \(A_{\mathfrak{p}}\) -模 \(M_{\mathfrak{p}}\) （对每个素理想 \(p \neq \mathfrak{m}_{\mathrm{A}}\)）是 Macaulay 的。 c) 设 M 为 Buchsbaum 模，\((x_{1}, \ldots, x_{n})\) 为对 M 割的一个序列；证明 \(\mathrm{M}/(x_{1}\mathrm{M} + \ldots + x_{n}\mathrm{M})\) 是一个 Buchsbaum 模。
\end{enumerate}

\(d\) ) 为使 M 是 Buchsbaum 模，必要且充分的是：对 M 的每个极大割序列 \((x_{1},\ldots ,x_{n})\) ，在 \(\mathrm{M}_n = \mathrm{M} / (x_1\mathrm{M} + \dots +x_{n - 1}\mathrm{M})\) 中有等式 \(\mathrm{Ker}(x_n)_{\mathrm{M}_n} = \mathrm{Ker}(x_n^2)_{\mathrm{M}_n}\) （先处理 \(n = 1\) 的情形，注意此时 \(\mathrm{Ker}(x_1)_{\mathrm{M}}\) 具有有限长度，并对这个长度用归纳法推理。在一般情形，把每个割序列 \((x_{1},\ldots ,x_{i})\) 补全为极大割序列 \((x_{1},\ldots ,x_{n})\) ，并考虑序列 \((x_{1},\ldots ,x_{i - 1},x_{i + 1}^{k},\ldots ,x_{n}^{k},x_{i})\) （对 \(k\geqslant 1\) ）。）

¶ 8) 设 A 为 Noether 局部环，M 为有限生成 A-模。我们打算证明下列两个条件的等价性：

\begin{enumerate}
\def\labelenumi{(\roman{enumi})}
\item
  M 是 Buchsbaum 模；
\item
  存在一个正整数 \(i(M)\) ，使得有 \(\lg(M/qM)-e_{\mathfrak{q}}(M)=i(M)\) （对每个理想 \(\mathfrak{q}\subset\mathfrak{m}_{A}\) ，它使 \(M/qM\) 具有有限长度）。
\end{enumerate}

\begin{enumerate}
\def\labelenumi{\alph{enumi})}
\item
  在条件 (i) 之下，设 x 为对 M 割的一个序列，并用 x 表示它所生成的理想。利用习题 6, d) 的记号，证明等式 \(\lg(\mathrm{M}/x\mathrm{M}) - e_{\mathfrak{x}}(\mathrm{M}) = \lg(\mathrm{K}_{n})\) 。设 \(\mathbf{x}'\) 为对 M 割的另一个序列，并设 \(\mathrm{K}_{n}^{\prime}\) 为相应的模；证明等式 \(\lg(\mathrm{K}_{n}) = \lg(\mathrm{K}_{n}^{\prime})\) （对 \(n\) 用归纳法推理，考虑一个元素 \(t\) （属于 \(\mathfrak{m}_{\mathrm{A}}\) ），使得序列 \((x_{1},\ldots,x_{n-1},t)\) 与 \((x_{1}',\ldots,x_{n-1}',t)\) 都对 M 割）。
\item
  我们假设条件 (ii) 成立。设 \(\mathbf{x} = (x_{1}, \ldots, x_{n})\) 为对 M 的一个极大割序列，并设 \(M_{n} = \mathrm{M}/(x_{1}\mathrm{M} + \ldots + x_{n-1}\mathrm{M})\) ；对序列 x 与 \((x_{1}, \ldots, x_{n-1}, x_{n}^{2})\) 应用习题 6, d) 的公式（并注意到习题 6, e)），证明有 \(\operatorname{Ker}(x_{n})_{\mathrm{M}_{n}} = \operatorname{Ker}(x_{n}^{2})_{\mathrm{M}_{n}}\) 。最后由习题 7, d) 得出结论）。
\item
  为使 M 是 Buchsbaum 模，必要且充分的是 \(\widehat{\mathbf{M}}\) 也是 Buchsbaum 模；此时我们有 \(i(\widehat{\mathbf{M}}) = i(\mathbf{M})\) 。
\item
  设 M 为 Buchsbaum 模，\((x_{1},\ldots,x_{p})\) 为 \(\mathfrak{m}_{A}\) 中元素组成的对 M 割的一个序列，且 \(p<\dim(M)\) ；证明我们有 \(i(\mathrm{M}/(x_{1}\mathrm{M}+\ldots+x_{p}\mathrm{M}))=i(\mathrm{M})\) （利用 a) 中给出的 \(i(\mathrm{M})\) 的表达式）。
\end{enumerate}

\begin{enumerate}
\def\labelenumi{\arabic{enumi})}
\setcounter{enumi}{8}
\tightlist
\item
  设 A 为 Noether 局部环。
\end{enumerate}

\begin{enumerate}
\def\labelenumi{\alph{enumi})}
\item
  设 M 为一个有限生成 Macaulay A-模，N 为 M 的一个包含 \(\mathfrak{m}_{A}M\) 的子模。证明 N 是一个 Buchsbaum 模，且有 \(i(N) = (\dim(M) - 1)[M/N : \kappa_{A}]\) 。
\item
  若 A 是 Macaulay 环，则理想 \(\mathfrak{m}_{A}\) 是一个 Buchsbaum 模，且当 \(\dim(A) \geqslant 2\) 时它不是 Macaulay 的。
\item
  设 \(k\) 为一个域；我们取 A 为 \(k[[X,Y]]\) 的由 \(X^4, X^3Y, XY^3, Y^4\) 生成的子环。有 \(\dim(A) = 2\) ，\(\text{prof}(A) = 1\) ，且 A 是一个 Buchsbaum 环，其 \(i(A) = 1\) （对 M 取 A 的整闭包并应用 \(a\) ）。
\end{enumerate}

¶ 10) a) 设 \(\rho: A \to B\) 为 Noether 环之间的一个单射同态，使得 B 的子 A-模 \(\mathrm{A}1_{B}\) 是一个直因子。证明若 B 是 Macaulay 环，则 A 也是。

（化为 A 是局部的情形，并通过对 \(\mathrm{prof}_{\mathrm{A}}(\mathrm{B})\) 用归纳法推理，化为 \(\mathrm{Ass}(\mathrm{B})\) 含有一个极大理想的情形。通过考虑 B 中 0 的准素分解，证明此时 B 同构于一个乘积，并对 B 的极大理想的个数用归纳法推理得出结论。

\begin{enumerate}
\def\labelenumi{\alph{enumi})}
\setcounter{enumi}{1}
\tightlist
\item
  设 B 为 Macaulay 环，G 为 B 的一个有限自同构群，其阶在 B 中可逆。证明 G 的不变量环是一个 Macaulay 环。
\end{enumerate}

¶ 11) 设 B 为一个 Noether 整闭环，G 为 B 的一个有限自同构群，A 为 G 的不变量环。

\begin{enumerate}
\def\labelenumi{\alph{enumi})}
\tightlist
\item
  我们假设：对 \(\mathrm{A}\) 的每个高为 1 的素理想 p 以及 B 的每个位于 p 之上的素理想 q，赋值 \(v_{\mathfrak{q}}\) 相对于 \(v_{\mathfrak{p}}\) 是非分歧的（VI, § 8, 第 1 节）。设 \(i: C(A) \to C(B)\) 为除子类群上的典范同态（VII, § 1, 第 10 节）；证明 \(i\) 的核同构于上同调群 \(H^{1}(G,B^{*})\) （A, X, 第 111 页；设 L 为 B 的分式域；考虑 Z \(^{(G)}\) -模的正合列
\end{enumerate}

\[
0 \rightarrow \mathrm{B} ^ {*} \rightarrow \mathrm{L} ^ {*} \rightarrow \mathrm{D} (\mathrm{B}) \rightarrow \mathrm{C} (\mathrm{B}) \rightarrow 0,
\]

并注意 \(\mathrm{H}^{0}(\mathrm{G},\mathrm{D}(\mathrm{B}))\) 等同于 \(\mathrm{D}(\mathrm{A})\) ，且 \(\mathrm{H}^{1}(\mathrm{G},\mathrm{L}^{*})\) 为零）。

\begin{enumerate}
\def\labelenumi{\alph{enumi})}
\setcounter{enumi}{1}
\item
  设 k 为一个特征为 2 的域；取 B 为环 \(k[X_{1},\ldots,X_{4}]\) ，并取 G 为 k-代数 B 的由自同构 \(\sigma\) （它满足 \(\sigma(X_{i})=X_{i+1}\) （对 \(1\leqslant i\leqslant3\) ）且 \(\sigma(X_{4})=X_{1}\) ）所生成的自同构群。由 a) 推出 G 的不变量环 A 是唯一分解的。
\item
  设 m 为 \(\mathrm{A}\) 上理想 \((\mathrm{X}_1, \ldots, \mathrm{X}_4)\) 的迹；令 \(s = \mathrm{X}_1 + \ldots + \mathrm{X}_4\) ，\(t = (\mathrm{X}_1 + \mathrm{X}_3)(\mathrm{X}_2 + \mathrm{X}_4)\) ，\(u = (\mathrm{X}_1 + \mathrm{X}_2)(\mathrm{X}_3 + \mathrm{X}_4)(\mathrm{X}_1 + \mathrm{X}_4)(\mathrm{X}_2 + \mathrm{X}_3)\) ，\(v = \mathrm{X}_1\mathrm{X}_2\mathrm{X}_3\mathrm{X}_4\) 。证明序列 \((s, t, u, v)\) 对 \(\mathrm{A}_{\mathfrak{m}}\) 是割的但不是完全割的。（设 \(x = \mathrm{X}_1^2(\mathrm{X}_2 + \mathrm{X}_3) + \mathrm{X}_2^2(\mathrm{X}_3 + \mathrm{X}_4) + \mathrm{X}_3^2(\mathrm{X}_4 + \mathrm{X}_1) + \mathrm{X}_4^2(\mathrm{X}_1 + \mathrm{X}_2)\) ，\(y = \mathrm{X}_1\mathrm{X}_3 + \mathrm{X}_2\mathrm{X}_4\) ，\(w = \mathrm{X}_1\mathrm{X}_2\mathrm{X}_3 + \mathrm{X}_2\mathrm{X}_3\mathrm{X}_4 + \mathrm{X}_3\mathrm{X}_4\mathrm{X}_1 + \mathrm{X}_4\mathrm{X}_1\mathrm{X}_2\) ；证明我们有 \(sw + y^2 = u\) ，且 \(xy + tw\) 被 \(s\) 整除，从而 \(ux\) 属于理想 \((s, t)\) 。）
\item
  证明 A 不是 Macaulay 环，且子 A-模 \(A.1_{B}\) 不是 B 的直和项。
\end{enumerate}

